\section{Simulator}
\label{sec:simulator-details}

The simulator runs the ledger and protocol of Sections~\ref{sec:ledger}
and~\ref{sec:construction}, except for the seven replacement hooks defined
in Section~\ref{sec:simulator-hooks}.  Figure~\ref{fig:simulator-patches}
shows where these hooks are used.  We first specify the logistics of
connecting this execution to $\FMPC$: registration, timing, input
commitment, and corruption disclosure.  These rules call the hooks;
all cryptographic replacements are defined in the hooks themselves.
The variants used in the proof replace those same named hooks.

\begin{figure}[p]
\centering
\begingroup
\definecolor{YasoControl}{HTML}{70516F}
\resizebox{\linewidth}{!}{%
\begin{tikzpicture}[x=1cm,y=1cm,font=\fontsize{9}{11}\selectfont,
 >={Latex[length=1.7mm]},
 entity/.style={draw=black!65,rounded corners=.8mm,line width=.6pt,fill=white},
 flow/.style={->,draw=YasoControl,line width=.75pt},
 lab/.style={align=left,text=YasoControl,fill=white,inner sep=2pt}]
\draw[entity] (0,0) rectangle (18,-.65);
\node at (9,-.325) {Environment $\mathcal Z$};
\draw[entity,draw=YasoControl] (0,-3.25) rectangle (18,-20.05);
\node[font=\sffamily\bfseries\large,text=YasoControl] at (9,-3.7)
 {Simulator $\mathcal S$: replacements in the protocol};
\draw[entity] (0,-23.65) rectangle (18,-24.35);
\node at (9,-24) {$\FMPC[G]$};

\draw[flow] (.35,-.65)--(.35,-3.25);
\draw[flow] (17.65,-3.25)--(17.65,-.65);
\node[lab,anchor=west] at (.65,-1.9)
 {$\Register(\Class),\ \Read,\ \Activate(\pk),\ \Post(\pk,m)$\\
  Corruption of $\pid$; random-oracle query $q$};
\node[lab,anchor=east,align=right] at (17.35,-1.9)
 {$\pk$, ledger $L$; oracle answer\\
  Exposed keys, tapes, inputs, and local histories};

\node[anchor=north west,inner sep=0pt,text width=16.8cm,align=left] at (.6,-4.25) {
Maintain the ledger's public state and execute its scheduling, posting,
selection, reconstruction, and evaluation rules (Figures~\ref{fig:yaso-ledger}, \ref{fig:role-allocation}--\ref{fig:protocol-steps}).
Mirror application registration and timing through $\FMPC$; commit honest
postings and submit extracted corrupt inputs.  $\FMPC$ delivers outputs
to honest pids; simulate queued deliveries to pids corrupted before their
delivery tick, as described in the text.
};
\draw[black!20] (.6,-5.9)--(17.4,-5.9);

\node[anchor=north west,inner sep=0pt,text width=16.8cm,align=left] at (.6,-6.15) {
\renewcommand{\arraystretch}{1.12}
\setlength{\tabcolsep}{0pt}
\begin{tabular}{@{}>{\raggedright\arraybackslash}p{3.4cm}>{\raggedright\arraybackslash}p{13.4cm}@{}}
$\Sampling_0$ & Use $\SimSetup$ and $\SimRO$ with a memoized sample source.
Before $Q_s$ is released, reserve the exact description $\SampleCircuit_s$
for the bundle prepared below.  Sample all other descriptions honestly.
The nonce-only sample is honest; retained nonce keys reveal $Q_s$ to $\mathcal S$.\\[7pt]

$\InputEncryption_0$ & Replace the input-key generation in Figure~\ref{fig:distribution-circuit} and
secret-input encryption in Figure~\ref{fig:signed-encrypted-inputs} by full-NCE simulation.  For an honest posting, sign
the public input together with the reserved ciphertext; explain its sender
tape on later corruption.  If corruption comes first, obtain a receiver key from a dummy-message
explanation and use it to interpret accepted adversarial ciphertexts.\\[7pt]

$\StateCiphertext_0$ & Generate state-transfer keys normally.
At input closure, choose fresh coins and set
$c_s=\PKE.\Enc(\statepk_{s+1},0^{|\sigma_s|})$.
Supply this value to $\Garbling_0$ as the state-ciphertext output
of Figure~\ref{fig:round-circuit}.\\[7pt]

$\OutputPads_0$ & At input closure, choose the masked bit $z_i$ uniformly
if its recipient's pad is not fixed.  On subsequent corruption, set
$r_i=z_i\oplus\ysec_i$.  If corruption comes before closure, fix a uniform
$r_i$ then use $z_i=\ysec_i\oplus r_i$ at closure.  Never revise a fixed
pad or masked bit.\\[7pt]

$\Garbling_0$ & In Figure~\ref{fig:distribution-circuit} use
$(\GC_s,\SimState_s)\sample\SimGarble(1^\lambda,\Topo(\RoundCircuit_s))$.
At input closure, call
$X_s\sample\SimEncode(\SimState_s,(z_s,c_s,\ypub_s))$ once.
Public reconstruction and $\Eval$ in Figure~\ref{fig:model-simulation}1 are unchanged.\\[7pt]

$\Sharing_0$ & During sample preparation fix uniform shares at the
initially corrupt positions.  Once $X_s$ is available, use
$\CompleteShares$ for each selected label and its common encoding material.
Pass these share packets to $\Deliveries_0$ before the selected
members speak.\\[7pt]

$\Deliveries_0$ & Replace delivery-key generation and the designated
label/pad encryptions in Figures~\ref{fig:conditional-release}, \ref{fig:distribution-circuit}, \ref{fig:protocol-steps} by RNCE simulation.  Reserve one ciphertext per
unexposed key and explain its packet before exposure.  Initially corrupt
keys are genuine.  Early exposure and nonce-only or unused roles explain
a dummy packet, discard its ciphertext, and use ordinary later encryption.
Application registration supplies its NCE helper through $\FMPC$.\\
\end{tabular}
};

\draw[black!20] (.6,-17.0)--(17.4,-17.0);
\node[anchor=north west,inner sep=0pt,text width=16.8cm,align=left] at (.6,-17.25) {
\textbf{Retained state.} Record input-arrival positions and public components;
obtain secret contents on corruption.  Save each explained key and tape once,
and reconstruct its prescribed local histories consistently on disclosure.
\par\smallskip
\textbf{Proof variants.} Restore input encryption, real state plaintexts,
early pads and state-encryption coins, real garbling, early sharing,
real deliveries, and finally real universal sampling.  Section~\ref{sec:hybrids} defines these replacements as
complete experiments $H_1,\ldots,H_7$.  Keep the specified aborts for threshold
failure, a previously answered designated description, and an accepted
new signed message before key exposure; correctness aborts enter in the proof.
};

\draw[flow] (.35,-20.05)--(.35,-23.65);
\draw[flow] (17.65,-23.65)--(17.65,-20.05);
\node[lab,anchor=west] at (.65,-21.85)
 {$\Register(),\ \Announce(s),\ \Proceed(s)$\\
  $\CommitInput(s,a)$\\$\CorruptInput(s,a,x)$; corruption of $\pid$};
\node[lab,anchor=east,align=right] at (17.35,-21.85)
 {$\pk,\widetilde c,\SimState_P,\sigsk,\rho_{\mathsf{sig}}$\\
  $(s,F_s,\boldsymbol{\pid}_s,\ReadyTick_s)$\\
  Input notices with $\xpub$; $\SelectedInput$ indices\\
  $(s,\xpub_s,\ypub_s)$; corrupt outputs and records};
\end{tikzpicture}}
\endgroup
\caption{The simulator of Figure~\ref{fig:model-simulation}, enlarged.
The named entries replace operations of the protocol; application I/O
and clock connections remain as in that figure.  The simulator uses only
the information supplied through its ideal interface.}
\label{fig:simulator-patches}
\end{figure}


\subsection{Execution logistics}
\label{sec:simulator-logistics}

Use the ledger's state and handlers from Figure~\ref{fig:yaso-ledger}.
All ideal-task commands below are sent to $\FMPC$, with session tag
$\sid$, through the interface of Figures~\ref{fig:mpc-registration}
and~\ref{fig:mpc-execution}.  Ledger commands and replies use the
adversarial interface in Figure~\ref{fig:model-simulation}.
In Figure~\ref{fig:construction-crs}, $\CRS$, replace $\Setup^H$
by $\Sampling_0.\Setup()$; use the resulting $\crs=(U)$ in
Figure~\ref{fig:yaso-ledger}, \emph{Initialization}.  Answer oracle queries through
$\Sampling_0.\QueryRO(q)$.  Retain input notices indexed by $(\pid,j)$,
selected input indexes, and the ideal results of completed rounds.
For each input role retain flags recording whether it posted honestly
and whether its input has committed, initially false.  Unknown input
contents and private outputs are retained symbolically until corruption
reveals them.

\begin{description}
\item[Registration.]
 Modify Figure~\ref{fig:yaso-ledger}, $\Register$, by replacing
 $\Gen$ with the delivery hook.  For a computation registration, sample
 its corruption flag $\beta$ with probability $p$ and call
 $\Deliveries_0.\Register(\Computation,\beta,\bot)$.
 For an application registration, set $\beta=0$, send $\Register()$ to $\FMPC$
 (Figure~\ref{fig:mpc-registration}, $\Register$) and pass its returned
 helper $h$ to $\Deliveries_0.\Register(\Application,0,h)$.
 Return a rejected ideal registration without inserting a ledger entry.
 Use the returned pid in the ledger's registration bookkeeping, retaining
 $\beta$ and setting $\HasSpoken=\bot$.  Return the pid to $\mathcal Z$;
 the application-pid output is supplied by $\FMPC$'s $\Register$ handler.
 For an initially corrupted computation pid, disclose the state specified
 by the ledger's $\Register$ handler.

\item[Input arrivals and selection.]
 Process the notices from Figure~\ref{fig:mpc-execution}, the two
 $\Input$ handlers, as follows.  On $(\Input,\pid,j,s,a,\xpub)$, store
 $(|L|,s,a,\xpub)$ at index $(\pid,j)$; on $(\Input,\pid,j,\bot)$ store
 $(|L|,\bot)$.  The ledger position is the arrival tag used in
 Figure~\ref{fig:yaso-ledger}, $\Input$.
 Store each selected index from $(\SelectedInput,s,a,j)$, sent by
 Figure~\ref{fig:mpc-execution}, \emph{Clock tick}, item~4.
 Forward to $\mathcal Z$ the inputs that those $\Input$ handlers route
 from already corrupted pids.  Honest input notices remain internal.

\item[Clock, allocation, and announcements.]
 Follow Figure~\ref{fig:yaso-ledger}, \emph{Clock tick} and $\EndRound$.
 When $\PublicData$ makes a new allocation of round $s$
 (Figure~\ref{fig:protocol-schedule}, item~2, calling
 Figure~\ref{fig:role-allocation}, $\Allocate$), send $\Announce(s)$
 and call $\Sampling_0.\Reserve(s)$.
 When $\CoinPhase(s+1)$ starts under $\PublicData$, item~3, send
 $\Proceed(s)$ once.  These ideal commands are defined in
 Figure~\ref{fig:mpc-registration}; they set
 $\ReadyTick_s=\AllocRound_s+1$ and
 $\InputStart_s=r^{\mathsf{coin}}_{s+1}+1$.
 Both functionalities follow the same clock.

\item[Honest activations.]
 Apply the eligibility guards of Figure~\ref{fig:yaso-ledger},
 $\Activate$, using the snapshots frozen by its \emph{Clock tick} handler.
 Computation parties run the release tests of
 Figure~\ref{fig:protocol-steps}, $\ComputationStep$, and return their
 cached delivery key when the test holds.  The sampling and sharing
 hooks fix that key before this release.
 In the input branch of $\ApplicationStep$ in the same figure, use
 the selected index received from $\FMPC$; remain silent if it is $\bot$.
 Otherwise replace $\Inputs.\SubmitInput$'s \emph{What} in
 Figure~\ref{fig:signed-encrypted-inputs} by
 $\InputEncryption_0.\Post(s,a,\xpub_{s,a})$ and post its signed record.
 Mark the input committed.  Send $\CommitInput(s,a)$ as defined in
 Figure~\ref{fig:mpc-execution} unless this activation is in the closing
 sweep; that sweep uses the automatic commitment in the same figure's
 \emph{Clock tick}, item~1.
 In the announcement branch of $\ApplicationStep$
 (Figure~\ref{fig:protocol-steps}), queue its circuit notification.
 For an assigned output recipient in $\OutputPhase(s)$ with the required
 public data available, replace private decryption in that figure's
 output branch by queuing the ideal output symbolically by role.  Apply the flags, postings, history updates, and output queue
 of Figure~\ref{fig:yaso-ledger}, $\Activate$; delivery is deferred to
 its $\EndRound$.  Other cases of $\ApplicationStep$ are unchanged.

\item[Adversarial postings.]
 Apply Figure~\ref{fig:yaso-ledger}, $\Post$.
 After an accepted post in $\InputPhase(s)$, select records as in
 Figure~\ref{fig:protocol-schedule}, $\PublicData$, item~1,
 $\InputPhase(s)$, but use $\PostedMessages(L,r)$ from the current
 ledger instead of the closed snapshot.  For each uncommitted provider
 $a$, use the context $\ctx_{s,a}$ and verification key $\sigvk_{s,a}$
 specified there in the call to $\Inputs.\Select$
 (Figure~\ref{fig:signed-encrypted-inputs}, \emph{Primitives}).
 Skip default results and apply the authentication check
 $\PO{3}$ in Section~\ref{sec:simulator-checks} to selected records.
 If the provider is corrupt, call
 $\InputEncryption_0.\Extract(s,a,c)$ and send
 $\CorruptInput(s,a,(\xsec,\xpub))$
 (Figure~\ref{fig:mpc-execution}) with the returned secret part and
 the selected public part.  Mark that provider committed.

\item[Input closure.]
 Retain the result and corrupt-recipient output bits supplied by
 Figure~\ref{fig:mpc-execution}, \emph{Clock tick}, item~1.
 Once this result and the closed ledger input records are available,
 form $\InputBits_s$ as in Figure~\ref{fig:protocol-schedule},
 $\PublicData$, item~1, $\InputPhase(s)$, and perform once:
 \[
 \begin{aligned}
 z_s&\gets\OutputPads_0.\Close(s,\boldsymbol y_s^{\mathsf{cor}}),\\
 c_s&\gets\StateCiphertext_0(s),\\
 X_s&\gets\Garbling_0.\Encode(s,(z_s,c_s,\ypub_s)),\\
 &\Sharing_0.\Open(s,\InputBits_s,X_s).
 \end{aligned}
 \]
 Test readiness after either prerequisite arrives and complete these
 calls before handling an honest label disclosure.  Recovery and
 evaluation remain as in $\PublicData$, item~1, $\LabelPhase(s)$.

\item[Corruption and private delivery.]
 Apply Figure~\ref{fig:yaso-ledger}, \emph{On corruption}, and mark
 the pid corrupted.  For an application pid, obtain its record through
 Figure~\ref{fig:mpc-execution}, \emph{On corruption}.  Merge its input
 contents with the stored arrival positions and store newly learned
 output bits in $\boldsymbol y_s^{\mathsf{cor}}$.
 For a prepared input provider call
 $\InputEncryption_0.\Corrupt(s,a,m)$, where $m$ is its committed secret
 input if it posted honestly, and $\bot$ otherwise.
 For an output recipient call $\OutputPads_0.\Corrupt(s,j,y)$, where
 $y$ is its determined ideal output, or $\bot$ if not determined.
 For an unassigned application pid or an input provider's unused
 delivery key, call $\Deliveries_0.\Bind(\pid,0^{\ell_{\mathsf{del}}})$.
 Reconstruct the state and history as in Section~\ref{sec:simulator-checks}
 and disclose them as specified by the ledger's corruption handler.
 Implement the delivery instruction of Figure~\ref{fig:yaso-ledger},
 $\EndRound$, using Figure~\ref{fig:mpc-execution}, \emph{Clock tick},
 items~2--3.  Keep the queue for slot $r-1$ through the deliveries at
 tick $r$.  $\FMPC$ delivers directly to honest pids.  On
 $(\Deliver,r,P,v)$ for a corrupt pid, check $\PendingOutputs_{r-1}$.
 If it contains an entry for $P$, consume it and immediately forward
 $(P,v)$ to $\mathcal Z$ as the ledger's output to $P$.
 Otherwise discard the notice.
\end{description}

Registration helpers, selected-input indexes, and early ideal results
are internal simulator state.  The timing of outputs to $\mathcal Z$
is that of the referenced ledger and ideal-task handlers.

\subsection{Replacement hooks}
\label{sec:simulator-hooks}

The following seven entries define the cryptographic replacements.
Their methods are local simulator procedures, called at the points in
Figure~\ref{fig:simulator-patches}.
A designated sample has exactly the five-component output tuple of
Figure~\ref{fig:distribution-circuit}.  State-transfer key generation,
signing-key generation, and the public algorithms not replaced below
are executed as in the protocol.

Write $\ncek_P$ for the RNCE key of registered pid $P$ (the first
component for an application pid).  A share packet uses the fixed padded
encoding of $\ESCR$; a pad packet is
$r\mathbin\|0^{\ell_{\mathsf{del}}-1}$.
Qualify the input and delivery NCE simulators by $\SNCE$ and $\RNCE$,
respectively.  Each explanation, ciphertext, and random tape is cached
once; repeated disclosures or sample requests reuse it.

\begin{description}
\item[$\Sampling_0$.]\phantomsection\label{hook:sampling}
 Replace public setup and oracle answers.  Maintain a sample table $T$
 indexed by exact descriptions, and a reservation table $R$.
 \begin{enumerate}
 \item Define $\SimSample(d)$: return $T[d]$ if defined.  Otherwise,
   draw uniform $\omega$ of the prescribed length.  If $R[d]=s$, execute
   $\SampleCircuit_s(\omega)$ with the replacements below; for any
   other description execute $d(\omega)$ with the preliminaries'
   invalid-description convention.
   Store the result in $T[d]$ and return it.
 \item $\Setup()$ runs
   $(U,\SimState_U)\sample\SimSetup^{\SimSample}(1^\lambda)$
   and returns $U$ for the ledger setup.  $\QueryRO(q)$ updates
   $\SimState_U$ and returns the answer from
   $\SimRO^{\SimSample}(\SimState_U,q)$.
 \item $\Reserve(s)$ takes the newly fixed allocation.  Check the
   nonce and label committees' thresholds.  For each initially honest
   nonce member $P$, call
   $\Deliveries_0.\Bind(P,0^{\ell_{\mathsf{del}}})$.
   Compute $\boldsymbol c\gets
   \Coin.\Prepare(U,\boldsymbol{\pk}^{\mathsf{coin}}_s)$ and
   $Q_s\gets\Coin.\Recover(\boldsymbol c,\boldsymbol{\sk}^{\mathsf{coin}}_s)$,
   using all retained nonce-member keys.  Form $d=\SampleCircuit_s$ from
   this nonce and the allocation.  Abort if $T[d]$ exists; otherwise
   set $R[d]=s$.  Finish before any honest nonce member speaks.
 \end{enumerate}

\item[$\InputEncryption_0$.]\phantomsection\label{hook:input-encryption}
 Replace sample-owned input-key generation and honest input submission;
 provide the key or sender tape needed on corruption.
 \begin{enumerate}
 \item $\Prepare(s,a)$ replaces input-key generation in
   $\SampleCircuit_s$ by
   $(\inputpk_{s,a},\widetilde c_{s,a},\tau_{s,a})
      \sample\SNCE.\SimNCE(1^\lambda)$ at length $\ell_{s,a}$.
   Retain these values and draw the prescribed independent decryption
   tape $\delta_{s,a}$, replacing its draw in step~2 as well.  If the provider is already corrupt, immediately
   perform the no-honest-posting branch of $\Corrupt$ below.  Return
   $\inputpk_{s,a}$.
 \item $\Post(s,a,\xpub)$ replaces the input interaction's ciphertext
   by $\widetilde c_{s,a}$.  Using the prescribed context
   $\ctx_{s,a}=(\sid,s,a,\SampleCircuit_s)$, form
   \[
     q=\InputRecord(\ctx_{s,a},\xpub,\widetilde c_{s,a}).
   \]
   Sign $q$ with the provider's registered signing key and fresh retained
   coins, record $q$ as honestly signed, and mark the provider as having
   posted honestly.  Return $(q,\sig)$.
 \item $\Corrupt(s,a,m)$: after an honest posting, obtain and cache
   the sender tape from $\SNCE.\ExplainNCE(\tau_{s,a},m)$, where $m$
   is the committed secret input supplied on corruption.  If no honest
   posting occurred, instead explain $0^{\ell_{s,a}}$, rerun
   $\SNCE.\KeyGen$ with its receiver tape, and retain
   $\inputsk_{s,a}$ internally.  Its reserved ciphertext is unused.
   The input-decryption key and tape are not input-provider state.
 \item $\Extract(s,a,c)$ uses this retained receiver key and
   $\delta_{s,a}$ to compute $\SNCE.\Dec(\inputsk_{s,a},c;\delta_{s,a})$;
   replace a result outside $\bits^{\ell_{s,a}}$ by $0^{\ell_{s,a}}$.
   Return the resulting secret input.
 \end{enumerate}

\item[$\StateCiphertext_0$.]\phantomsection\label{hook:state-ciphertext}
 Defer outgoing state-encryption coins until input closure.
 The call $\StateCiphertext_0(s)$ returns
 \[
   c_s\sample\PKE.\Enc(\statepk_{s+1},0^{\nu_s}).
 \]
 This is the state-ciphertext component supplied to the garbling hook.
 State-transfer key generation in $\SampleCircuit$ is unchanged.

\item[$\OutputPads_0$.]\phantomsection\label{hook:output-pads}
 Replace pad generation and masked-output calculation.  Retain
 $r_{s,j},z_{s,j}$, initially unset, for each recipient.
 \begin{enumerate}
 \item $\Prepare(s)$ runs once at the pad draw in step~2 of
   $\SampleCircuit_s$.  Cache its returned vector as the output packets
   used in step~4.
   For each recipient $P=\pid^{\mathsf{out}}_{s,j}$ whose delivery
   key is already fixed, choose $r_{s,j}\sample\bits$ and set $m$ to
   its pad packet.  Otherwise leave the pad unset and set $m=\bot$.
   Set $\OutputPackets_{s,j}\gets\Deliveries_0.\Packet(P,m)$.
   Return the resulting output-packet vector.
 \item $\Close(s,\boldsymbol y^{\mathsf{cor}}_s)$ sets each unset
   $z_{s,j}$ uniformly if its pad is unset; otherwise set
   $z_{s,j}=y^{\mathsf{cor}}_{s,j}\oplus r_{s,j}$.
   In this latter case the recipient has already been corrupted.
   Return $z_s$.
 \item $\Corrupt(s,j,y)$: if no designated packet has been fixed,
   call $\Deliveries_0.\Bind(\pid^{\mathsf{out}}_{s,j},
   0^{\ell_{\mathsf{del}}})$ and return.  Otherwise, if the pad is unset,
   choose it uniformly when $z_{s,j}$ is unset, and set
   $r_{s,j}=z_{s,j}\oplus y$ when $z_{s,j}$ is fixed.  In the latter
   case $y$ is supplied by the ideal corruption record.  Call
   $\Deliveries_0.\Bind(\pid^{\mathsf{out}}_{s,j},
   r_{s,j}\mathbin\|0^{\ell_{\mathsf{del}}-1})$.
 \end{enumerate}

\item[$\Garbling_0$.]\phantomsection\label{hook:garbling}
 Replace the garbling and input encoding.
 The call $\Prepare(s)$ returns $\GC_s$ and retains $\tau_s$, using
 $(\GC_s,\tau_s)\sample\SimGarble(1^\lambda,\Topo(\RoundCircuit_s))$,
 computing the topology from the public dimensions and Boolean template.
 No secret constants need to be instantiated.
 The one call $\Encode(s,(z_s,c_s,\ypub_s))$ computes and retains
 \[
   X_s\sample\SimEncode(\tau_s,(z_s,c_s,\ypub_s)).
 \]
 Retain this prescribed public result and return $X_s$.

\item[$\Sharing_0$.]\phantomsection\label{hook:sharing}
 Replace $\Labels.\Prepare$ and the later determination of selected
 share packets.  Let $I_{s,w,b}$ index the initially corrupted members.
 \begin{enumerate}
 \item $\Prepare(s)$ chooses independent uniform Shamir shares
   $(u_i)_{i\in I_{s,w,b}}$ for each $(w,b)$.
   For member $P=\pk_{s,w,b,i}$, let $m$ be its encoded share if
   $i\in I_{s,w,b}$, and $\bot$ otherwise.  Set its entry of
   $\LabelPackets_s$ to $\Deliveries_0.\Packet(P,m)$.
   Return this array and retain the chosen shares.
 \item $\Open(s,x,X_s)$ parses
   $X_s=(\GarbleAux,\widehat L_1,\ldots,\widehat L_{m_s})$.
   For each $w$, put $b=x_w$ and sample
   \[
    (v_i)_{i\in[n]}\sample
      \CompleteShares(I_{s,w,b},(u_i)_{i\in I_{s,w,b}},
                       \langle\GarbleAux,\widehat L_w\rangle).
   \]
   Retain this selected sharing.  For each initially honest member
   $P=\pk_{s,w,b,i}$, call $\Deliveries_0.\Bind(P,m_i)$ with $m_i$
   the encoded share $v_i$, before it speaks.
   The unselected honest members remain unexplained and silent.
 \end{enumerate}

\item[$\Deliveries_0$.]\phantomsection\label{hook:deliveries}
 Replace registered delivery keys and designated label/pad ciphertexts.
 For every pid $P$, retain a receiver key and tape, initially unset for
 simulated keys, and an initially unset designated ciphertext $D_P$.
 \begin{enumerate}
 \item $\Register(\Class,\beta,h)$ returns the registered pid.
   Here $\beta\in\bits$ is the computation-corruption flag and $h=\bot$
   for computation registrations; for application registrations,
   $\beta=0$ and
   $h=(P,\widetilde c_P,\tau_P,\sigsk_P,\rho^{\mathsf{sig}}_P)$ is the
   helper returned by $\FMPC$.  For a computation pid with $\beta=1$ use
   genuine $\RNCE.\KeyGen$ and retain its key and complete tape.
   For a computation pid with $\beta=0$ use
   $(\ncek_P,\widetilde c_P,\tau_P)\sample\RNCE.\SimNCE(1^\lambda)$.
   For an application pid use exactly the pid and helper supplied by
   ideal registration.  Retain its genuine signing material separately.
 \item $\Packet(P,m)$, with $m\in\bits^{\ell_{\mathsf{del}}}\cup\{\bot\}$:
   if $D_P$ exists, return it.  If the receiver key is unset, set
   $D_P=\widetilde c_P$.  Otherwise require a known $m$ and set
   $D_P\sample\RNCE.\Enc(\ncek_P,m)$.  Return $D_P$.
   The ciphertext is fixed when included in the memoized sample,
   independently of when that sample is first observed publicly.
 \item $\Bind(P,m)$, with $m\in\bits^{\ell_{\mathsf{del}}}$:
   if the receiver key is unset, call
   $(\rho_R,\rho_S)\sample\RNCE.\ExplainNCE(\tau_P,m)$,
   rerun $\RNCE.\KeyGen(1^\lambda;\rho_R)$, and retain its secret key
   and complete receiver tape.  Return the cached key.
   A dummy binding discards the unused reserved ciphertext; subsequent
   nonce ciphertexts or a later designated packet use ordinary encryption.
 \end{enumerate}
\end{description}

\subsection{Checks and retained histories}
\label{sec:simulator-checks}

The following checks define proof obligations $\PO{1}$--$\PO{9}$.
The simulator halts at their first failure; the guarded hybrid experiments
return bit $0$.  Section~\ref{sec:hybrids-details} bounds their probabilities in the
corresponding complete experiments.
\begin{description}
 \item[$\PO{1}$: committee thresholds.]\phantomsection\label{po:threshold}
 At reservation, check that each nonce and label committee has at most
 $t$ initially corrupted members.  These are the blocks assigned by
 Figure~\ref{fig:role-allocation}, $\Allocate$, item~3.
 \item[$\PO{2}$: freshness.]\phantomsection\label{po:freshness}
 In $\Sampling_0.\Reserve(s)$, check that $T[\SampleCircuit_s]$ is unset.
 The table includes source calls made by $\SimSetup$ and $\SimRO$.
 \item[$\PO{3}$: authentication.]\phantomsection\label{po:authentication}
 When $\Inputs.\Select$ accepts a record under an unexposed signing
 key, check that its message was signed by the simulated honest posting
 step.  The selection rule is in Figure~\ref{fig:signed-encrypted-inputs},
 \emph{Primitives}, $\Inputs.\Select$.  A different signature on the
 same signed message does not fail this check.
 \item[$\PO{4}$: public samples.]\phantomsection\label{po:samples}
 At every sampling call made by the simulated ledger or an honest step,
 compare its public result with $\SimSample(d)$.  Include the private
 $\Coin.\Prepare$ evaluation in $\Sampling_0.\Reserve$.  The public
 algorithms are Figure~\ref{fig:shared-coins}, $\Coin.\Prepare$, and
 Figure~\ref{fig:protocol-schedule}, $\PublicData$, item~1,
 $\CoinPhase(s)$.
 \item[$\PO{5}$: nonce recovery.]\phantomsection\label{po:nonce}
 At completion of $\CoinPhase(s)$ in $\PublicData$, item~1, check
 its recovered nonce against the reserved $Q_s$.
 \item[$\PO{6}$: input selection.]\phantomsection\label{po:inputs}
 At input closure, check that every honest posting is the selected
 record for its provider, and that every still-honest provider with a
 selected input posted.  Compare the public-input vector from
 Figure~\ref{fig:protocol-schedule}, $\PublicData$, item~1,
 $\InputPhase(s)$, with that supplied by Figure~\ref{fig:mpc-execution},
 \emph{Clock tick}, item~1.
 \item[$\PO{7}$: label recovery and evaluation.]\phantomsection\label{po:evaluation}
 At completion of $\LabelPhase(s)$ in Figure~\ref{fig:protocol-schedule},
 $\PublicData$, item~1, compare the reconstructed encoding with the
 retained $X_s$, and its evaluation with the garbling hook's retained
 result.
 \item[$\PO{8}$: announcements.]\phantomsection\label{po:announcements}
 Compare the circuit and pid vectors from
 Figure~\ref{fig:mpc-registration}, $\Announce$, with the simulated
 allocation, and check $\ReadyTick_s=\AllocRound_s+1$.
 This is the delivery tick for the notification branch of
 Figure~\ref{fig:protocol-steps}, $\ApplicationStep$, under
 Figure~\ref{fig:yaso-ledger}, $\EndRound$.
 \item[$\PO{9}$: disclosed histories.]\phantomsection\label{po:histories}
 Before corruption disclosure, check the reconstructed histories against
 their fixed public keys, postings, and queued outputs, as specified by
 Figure~\ref{fig:yaso-ledger}, $\Register$, $\Activate$, and
 \emph{On corruption}.
\end{description}
Ill-formed or unavailable required data fails its corresponding check.
An expected ideal notification is awaited; readiness is tested after
either its arrival or the corresponding ledger event.  No check requires
an honest provider's unrevealed sample-owned input-decryption key.

To reconstruct the histories required by Figure~\ref{fig:yaso-ledger},
$\Register$ and $\Activate$, retain each activation's ledger snapshot,
input indexes, used tape positions, and public-oracle query/answer
transcript.
On corruption, fill the missing input contents from the ideal record
(Figure~\ref{fig:mpc-execution}, \emph{On corruption}) and insert the
hook's explained coins into the corresponding subtapes from
Figure~\ref{fig:protocol-steps}.  Preserve all fixed positions and sample
unused positions independently.  Rerun $\Gen$ and the saved party steps
with these tapes and recorded oracle answers to obtain their local
configurations.  Disclose the
state required by the ledger's corruption handler, including only the
outputs computed in those activations.  As in Figure~\ref{fig:protocol-steps},
$\ComputationStep$, a computation member runs its release test and
returns its key; share completion and binding are simulator operations.
